%Este trabalho está licenciado sob a Licença Creative Commons Atribuição-CompartilhaIgual 3.0 Não Adaptada. Para ver uma cópia desta licença, visite https://creativecommons.org/licenses/by-sa/3.0/ ou envie uma carta para Creative Commons, PO Box 1866, Mountain View, CA 94042, USA.

\chapter{Introdução}

Cálculo numérico é a disciplina que estuda as técnicas para a solução aproximada de problemas matemáticos. Estas técnicas são de natureza analítica e computacional. As principais preocupações normalmente envolvem exatidão e desempenho.

Aliado ao aumento contínuo da capacidade de computação disponível, o desenvolvimento de métodos numéricos tornou a simulação computacional\index{simulação!computacional} de problemas matemáticos uma prática usual nas mais diversas áreas científicas e tecnológicas. As então chamadas simulações numéricas\index{simulação!numérica} são constituídas de um arranjo de vários esquemas numéricos dedicados a resolver problemas específicos como, por exemplo: resolver equações algébricas, resolver sistemas de equações lineares, interpolar e ajustar pontos, calcular derivadas e integrais, resolver equações diferenciais ordinárias etc. Neste livro, abordamos o desenvolvimento, a implementação, a utilização e os aspectos teóricos de métodos numéricos para a resolução desses problemas.

Trabalharemos com problemas que abordam aspectos teóricos e de utilização dos métodos estudados, bem como com problemas de interesse na engenharia, na física e na matemática aplicada.

A necessidade de aplicar aproximações numéricas decorre do fato de que esses problemas podem se mostrar intratáveis se dispomos apenas de meios puramente analíticos, como aqueles estudados nos cursos de cálculo e álgebra linear. Por exemplo, o teorema de Abel-Ruffini nos garante que não existe uma fórmula algébrica, isto é, envolvendo apenas operações aritméticas e radicais, para calcular as raízes de uma equação polinomial de qualquer grau, mas apenas casos particulares:
\begin{itemize}
\item Simplesmente isolar a incógnita para encontrar a raiz de uma equação do primeiro grau;
\item Fórmula de Bhaskara para encontrar raízes de uma equação do segundo grau;
\item Fórmula de Cardano para encontrar raízes de uma equação do terceiro grau;
\item Existe expressão para equações de quarto grau;
\item Casos simplificados de equações de grau maior que 4 onde alguns coeficientes são nulos também podem ser resolvidos.
\end{itemize}
Equações não polinomiais podem ser ainda mais complicadas de resolver exatamente, por exemplo:
\begin{equation}
\cos(x)=x\qquad \text{ou}\qquad xe^x= 10
\end{equation}

Para resolver o problema de valor inicial
\begin{equation}
\begin{array}{l}
y'+xy=x,\\
y(0)=2,
\end{array}
\end{equation}
podemos usar o método de fator integrante e obtemos $y=1+e^{-x^2/2}$. No entanto, não parece possível  encontrar uma expressão fechada em termos de funções elementares para a solução exata do problema de valor inicial dado por
\begin{equation}
\begin{array}{l}
y'+xy=e^{-y},\\
y(0)=2,
\end{array}.
\end{equation}


Da mesma forma, resolvemos a integral
\begin{equation}
\int_1^2xe^{-x^2}dx
\end{equation}
pelo método da substituição e obtemos $\frac{1}{2}(e^{-1}-e^{-4})$. Porém a integral
\begin{equation}
\int_1^2 e^{-x^2} dx
\end{equation}
não pode ser expressa analiticamente em termos de funções elementares, como uma consequência do teorema de Liouville.

A maioria dos problemas envolvendo fenômenos reais produzem modelos matemáticos cuja solução analítica é difícil (ou impossível) de obter, mesmo quando provamos que a solução existe. Nesse curso propomos calcular aproximações numéricas para esses problemas, que apesar de, em geral, serem diferentes da solução exata, mostraremos que elas podem ser bem próximas.

Para entender a construção de aproximações é necessário estudar como funciona a aritmética implementada nos computadores e erros de arredondamento. Como computadores, em geral, usam uma base binária para representar números, começaremos falando em mudança de base.


\section{A presença dos computadores}
O surgimentos dos computadores eletrônicos representou um grande empuxo para o desenvolvimento de métodos numéricos e permitiu resolver problemas de grande escala, como previsão do tempo, simulação do escoamento de fluidos, previsão de terremotos, projeto e simulação de reatores nucleares, a análise de estruturas, a simulação de sistemas biológicos, a previsão do comportamento de mercados financeiros, entre outros.

A primeira linguagem de programação de alto nível, o Fortran, foi desenvolvida na década de 1950 por John Backus e sua equipe na IBM. O Fortran foi lançado em 1957 e é considerada a primeira linguagem de programação a ser amplamente adotada e é ainda hoje uma das linguagens mais usadas em computação científica. Foi a primeira linguagem a ser padronizada, em 1966, e desde então passou por várias atualizações, todas definidas por comitês de padronização.


Por ter uma sintaxe muito simples e um conjunto de palavras-chave muito limitado, o Fortran permitia importantes optimizações de código para produzir programas eficientes em tempo de execução. Programas escritos em Fortran podiam ser significaticamente mais velozes que equivalentes escritos noutras linguagens.

A partir dos anos 1970, enquanto o Fortran perdia espaço para outros tipos de linguagens de programação fora dos meio acadêmicos, ele continuou a ser amplamente usado em computação científica e até meados dos anos 1980, se mantinha quase como a única escolha para computação científica. 

O Matlab foi desenvolvido no final dos anos 1970 por Cleve Moler, foi primeiramente distribuido gratuitamente para estudantes e professores como uma ferramenta de ensino. Depois for remodelado e transformado num produto comercial. O Matlab foi desenvolvido para ser uma linguagem de programação de alto nível, interpretada, que permite a manipulação de matrizes e vetores de forma muito eficiente. O Matlab segue muito usado em algumas áreas em especial na teoria de controle, tratamento de sinais e comunicações e processamento de imagens.

A linguagem C foi desenvolvida por Dennis Ritchie na década de 1970, e foi uma das primeiras linguagens de programação a ser usada para desenvolver sistemas operacionais. A linguagem C é uma linguagem de programação que permite um controle preciso sobre o hardware e é muito usada em sistemas embarcados, sistemas operacionais, drivers de dispositivos e softwares de base em geral. A linguagem C é frequentemente denominda de "língua franca" da programação porque existem compiladores para quase todas as arquiteturas de computadores e sistemas operacionais. A maioria das linguagens de programação são capazes de chamar funções escritas em C, o que faz com que a linguagem C seja muito usada para escrever bibliotecas de funções que podem ser chamadas por outras linguagens. 

Devido à sua portabilidade, popularidade e capacidade de produzir código altamente otimizado, a linguagem C se tornou muito usada em computação científica. Além de novos projetos que surgiam nativamente na linguagem, muitas bases de códigos de software científico anteriormente escritas em Fortran passaram a reescritas em C. O movimento foi acelerado pelas melhorias na linguagem e nas limitações do Fortran 77 (lançado em 1977), que não permitia a alocação dinâmica de memória, o uso de funções recursivas nem orientação a objetos. Quando uma nova versão do Fortran foi publicada em 1990, muitos equipes de pesquisa e desenvolvimento já haviam migrado para C e não voltaram para Fortran.

Ainda nos anos 1980, linguagens como Basic e Pascal ganharam popularidade, mas deixaram de ser usadas em computação científica e basicamente cairam em desuso. O Basic era uma linguagem de programação de alto nível, interpretada, que foi muito usada em microcomputadores pessoais, mas não era adequada para computação científica. O Pascal era uma linguagem de programação de alto nível, compilada, que foi muito usada em cursos de introdução à programação de computadores.

Em meados dos anos 1980, Bjarne Stroustrup desenvolveu a linguagem C++, que é uma extensão da linguagem C. A linguagem C++ é uma linguagem de programação de alto nível, compilada, que permite a programação orientada a objetos em outros paradigmas. Embora não tenha sido a pioneiras nesses paradigmas, a linguagem C++ foi a primeira linguagem mainstream a trazer orientação a objetos e programação genérica numa linguagem de tipagem estática.

O orientação a objetos foi inicialmente vista com ceticismo por muitos programadores em especial na comunidade científica, que temiam por perda de performance e por complexidade. No entanto, a orientação a objetos se mostrou muito útil representar estruturas de dados mais complexes e para organizar e reutilizar código. A adoção do C+ pelo CERN para escrever sua biblitoeca ROOT incentivou a comunidade de física de alta energia a migrar para a linguagem, o que ajudou a popularizar ainda mais a linguagem. O C++ foi a primeira linguagem de programação mainstream de compilação estática com suporte a escrever código genérico, que pode ser usado com diferentes tipos de dados. A linguagem C++ acabou se tornando muito popular em computação científica, em especial para desenvolver bibliotecas de funções que podem ser chamadas por outras linguagens.

O Python foi desenvolvido por Guido van Rossum no final dos anos 1980 e lançado em 1991. O Python é uma linguagem de programação de alto nível, interpretada, que é muito usada em computação científica, em especial em áreas como aprendizado de máquina, inteligência artificial, análise de dados, visualização de dados e computação simbólica. É uma linguagem de programação muito versátil, que permite a programação orientada a objetos, programação funcional e programação imperativa. 
% https://docs.python.org/3/howto/functional.html

A linguagem é muito usada em cursos de introdução à programação de computadores por sua popularidade e por sua sintaxe muito clara e concisa, que prioriza a legibilidade do código e expõe gradualmente as complexidades da programação. No entanto, o Python é uma linguagem multiparadigma, permite diversos estilos de programação, e é muito usada em projetos de grande escala. Influenciada também por linguagens de maior sofitisticação matemática como Lisp e Haskell, existem tipos nativos para representar diversas abstrações matemáticas como conjuntos, dicionários, sequêcias com semântica própria, isto é, permitindo manipulação via operadores lógicos como "para todo" e "existe".

Outras linguagens como o Scilab e Octave foram desenvolvidas como alternativas livres ao Matlab. O Scilab foi lançado nos anos 1990. O Octave foi desenvolvido no final dos anos 1980. A linguagem Julia foi lançada em 2012 e se propõe trazer uma sintaxe funcional como a do a Python e a eficiência do C. 

Atualmente, muitas aplicações são desenvolvidas em várias linguagens, a computação bruta, onde performace é crítica, é programada em C ou C++, enquanto a interface com o usuário é programada em Python ou outra linguagem de alto nível. A interoperacionamente das linguagens é bastante comum. Por exemplo, a solução de um problema matemático descrito em Python é tipicamente executada por um interpretador escrito em C, o CPython, que chama funções da biblioteca Laplack, escrita em Fortran, que por sua vez chama funções de uma biblioteca BLAS, tipicamente escrita em C e Assembly.

%\end{document}
